\chapter{Product Measures and Change of Variables}

\begin{center}\itshape
Cut a loaf of bread into thin slices and add up the areas of the slices: you get its volume. Cut it lengthwise instead, and you
must get the same number. Why, and what happens to volume when the loaf is squeezed or twisted?
\end{center}

In 1906 the Danish philologist Johan Ludvig Heiberg examined a prayer book in a library in Constantinople. Under the
prayers, faintly, he could read an older text that a thirteenth-century monk had scraped off to reuse the parchment: it was
Archimedes, and among the works was one thought lost forever, \emph{The Method}. In it Archimedes explains how he
\emph{found} the volume of the sphere before proving it: he cut the sphere, a cone and a cylinder into thin parallel slices
and balanced the slices against each other on an imaginary lever. Volumes by slices.

The following year, 1907, Guido Fubini published the theorem that makes slicing rigorous: the integral over a product is
an iterated integral, in either order. Leonida Tonelli added in 1909 that for non-negative functions no condition at all is
needed. Between Archimedes and Fubini stands Bonaventura Cavalieri, whose principle (1635) says that solids with slices of
equal areas have equal volumes; his ``indivisibles'' were attacked by the Jesuits as a dangerous idea, and the quarrel
lasted decades. Measure theory settles it peacefully: the measure of a set in a product is the integral of the measures of
its slices.

Tome~0, Chapter~5, swapped double series; Tome~3 admitted Tonelli and Fubini as (A2); Tome~4 used Tonelli for Ritelli's
proof of $\zeta(2)=\frac{\pi^2}6$ and for $\Gamma(x)\Gamma(y)=B(x,y)\Gamma(x+y)$. All of this is proved here. The second half of
the chapter asks what happens to volume when space is squeezed or twisted. A linear map multiplies volumes by $|\det T|$,
which is what a determinant really is, and a smooth bijection multiplies them, near each point, by its Jacobian. Out come
polar coordinates in every dimension, the volume of the $n$-ball, which tends to zero as the dimension grows, and the density
of a Gaussian vector, as Tome~3, Chapter~6, announced.

\section{Product \texorpdfstring{$\sigma$}{sigma}-algebras and sections}

Let $(E,\mathcal A)$ and $(F,\mathcal B)$ be measurable spaces. A \emph{measurable rectangle} is a set $A\times B$ with $A\in\mathcal A$, $B\in\mathcal B$. The \emph{product $\sigma$-algebra}
$\mathcal A\otimes\mathcal B$ is the $\sigma$-algebra on $E\times F$ generated by the measurable rectangles. Rectangles form a $\pi$-system, since $(A\times B)\cap(A'\times B')=(A\cap A')\times(B\cap B')$.
The projections $(x,y)\mapsto x$ and $(x,y)\mapsto y$ are measurable, and $\mathcal A\otimes\mathcal B$ is the smallest $\sigma$-algebra with this property, exactly as the product topology
was the coarsest topology making the projections continuous (Chapter~3).

\begin{proposition}\label{prop:borelproduct}
$\mathcal B(\R^m)\otimes\mathcal B(\R^n)=\mathcal B(\R^{m+n})$.
\end{proposition}
\begin{proof}
The rectangles $A\times\R^n$ with $A$ open, and $\R^m\times B$ with $B$ open, are open in $\R^{m+n}$; the sets $A$ for which $A\times\R^n$ is Borel form a $\sigma$-algebra containing the open sets, so
$A\times\R^n$ is Borel for every Borel $A$, and similarly $\R^m\times B$. Their intersection $A\times B$ is Borel: $\mathcal B(\R^m)\otimes\mathcal B(\R^n)\subseteq\mathcal B(\R^{m+n})$. Conversely, every open set of
$\R^{m+n}$ is a countable union of open boxes with rational corners, and each such box is a product of an open box of $\R^m$ and an open box of $\R^n$. So the open sets, hence the Borel
sets, lie in the product $\sigma$-algebra.
\end{proof}

For $Q\subseteq E\times F$, $x\in E$ and $y\in F$, the \emph{sections} are $Q_x=\{y\in F:(x,y)\in Q\}$ and $Q^y=\{x\in E:(x,y)\in Q\}$. For a function $f$ on $E\times F$, the sections are
$f_x=f(x,\cdot)$ and $f^y=f(\cdot,y)$.

\begin{lemma}\label{lem:sections}
If $Q\in\mathcal A\otimes\mathcal B$, then $Q_x\in\mathcal B$ and $Q^y\in\mathcal A$ for all $x,y$. If $f$ is $\mathcal A\otimes\mathcal B$-measurable, then $f_x$ and $f^y$ are measurable.
\end{lemma}
\begin{proof}
Fix $x$. The family of $Q$ with $Q_x\in\mathcal B$ contains every rectangle ($(A\times B)_x$ is $B$ or $\varnothing$) and is a $\sigma$-algebra, since taking sections commutes with complements
and countable unions. So it contains $\mathcal A\otimes\mathcal B$. For functions, $\{f_x\le t\}=\{f\le t\}_x$.
\end{proof}

This ``good sets'' argument will be used again and again: to prove a property of all measurable sets, show that the sets having it form a
$\sigma$-algebra (or a $\lambda$-system) and contain a generating family (or a generating $\pi$-system).

\section{Product measures}

A measure $\mu$ on $(E,\mathcal A)$ is \emph{$\sigma$-finite} if $E$ is a countable union of sets of finite measure; we may take them increasing. Lebesgue measure, probabilities, and
counting measure on a countable set are $\sigma$-finite; counting measure on $\R$ is not.

\begin{lemma}\label{lem:slicemeasurable}
Let $\mu$ and $\nu$ be $\sigma$-finite measures on $(E,\mathcal A)$ and $(F,\mathcal B)$, and $Q\in\mathcal A\otimes\mathcal B$. Then $x\mapsto\nu(Q_x)$ is $\mathcal A$-measurable, $y\mapsto\mu(Q^y)$ is $\mathcal B$-measurable, and
\[\int_E\nu(Q_x)\,d\mu(x)=\int_F\mu(Q^y)\,d\nu(y).\]
\end{lemma}
\begin{proof}
\emph{Finite measures.} Let $\mathcal D$ be the family of $Q\in\mathcal A\otimes\mathcal B$ for which both functions are measurable and the two integrals agree. For $Q=A\times B$, $\nu(Q_x)=\nu(B)\mathbf 1_A(x)$
and $\mu(Q^y)=\mu(A)\mathbf 1_B(y)$, and both integrals equal $\mu(A)\nu(B)$; so $\mathcal D$ contains the rectangles, including $E\times F$. If $Q\subseteq Q'$ are in $\mathcal D$, then $\nu((Q'\setminus Q)_x)=\nu(Q'_x)-\nu(Q_x)$,
a difference of finite measurable functions, and the integrals subtract (all are finite): $Q'\setminus Q\in\mathcal D$. If $Q_n\in\mathcal D$ increase to $Q$, then $\nu((Q_n)_x)$ increases to $\nu(Q_x)$, so $\nu(Q_x)$ is
measurable as a limit, and monotone convergence passes to the limit in both integrals: $Q\in\mathcal D$. So $\mathcal D$ is a $\lambda$-system containing the $\pi$-system of rectangles, and by
Dynkin's theorem (Theorem~\ref{thm:dynkin6}) it is all of $\mathcal A\otimes\mathcal B$.

\emph{$\sigma$-finite measures.} Let $E_n\uparrow E$ and $F_n\uparrow F$ with $\mu(E_n),\nu(F_n)<\infty$. Apply the finite case to the finite measures $\mu_n=\mu(\cdot\cap E_n)$ and $\nu_n=\nu(\cdot\cap F_n)$: the
functions $x\mapsto\nu_n(Q_x)\mathbf 1_{E_n}(x)=\nu((Q\cap(E_n\times F_n))_x)$ are measurable and increase to $\nu(Q_x)$, and similarly in $y$; monotone convergence finishes.
\end{proof}

\begin{theorem}[product measure]\label{thm:product}
Let $\mu$ and $\nu$ be $\sigma$-finite. There is a unique measure $\mu\otimes\nu$ on $\mathcal A\otimes\mathcal B$ with $(\mu\otimes\nu)(A\times B)=\mu(A)\nu(B)$ for all measurable rectangles (with $0\cdot\infty=0$). For every
$Q\in\mathcal A\otimes\mathcal B$,
\[(\mu\otimes\nu)(Q)=\int_E\nu(Q_x)\,d\mu(x)=\int_F\mu(Q^y)\,d\nu(y).\]
\end{theorem}
\begin{proof}
Define $(\mu\otimes\nu)(Q)$ by the first integral; Lemma~\ref{lem:slicemeasurable} says it equals the second. It vanishes on $\varnothing$, and if $Q=\bigsqcup_kQ_k$, then $Q_x=\bigsqcup_k(Q_k)_x$, so
$\nu(Q_x)=\sum_k\nu((Q_k)_x)$, and the series can be integrated term by term (monotone convergence): $\mu\otimes\nu$ is a measure. Its value on rectangles was computed in the lemma. Uniqueness: two such
measures agree on the $\pi$-system of rectangles, and $E\times F$ is the increasing union of the rectangles $E_n\times F_n$ of finite measure; apply Theorem~\ref{thm:uniq6}.
\end{proof}

\begin{figure}[ht]
\centering
\fig{ch08-cavalieri}
\caption{Cavalieri's principle. Each horizontal slice of the two solids has the same area, so they have the same volume:
$\lambda_3(Q)=\int\lambda_2(Q^t)\,dt$. A slanted stack of coins and a straight one hold the same amount of metal.}
\label{fig:cavalieri}
\end{figure}

\emph{Lebesgue measure as a product.} The measures $\lambda_{m+n}$ and $\lambda_m\otimes\lambda_n$ live on the same $\sigma$-algebra (Proposition~\ref{prop:borelproduct}) and agree on boxes, since the volume of a box is the
product of the volumes of its two factors. By uniqueness (Theorem~\ref{thm:lebesgue}), $\lambda_{m+n}=\lambda_m\otimes\lambda_n$, and by induction $\lambda_d=\lambda\otimes\cdots\otimes\lambda$. This is the second
half of Tome~3's (A1). Theorem~\ref{thm:product} becomes \emph{Cavalieri's principle} (Figure~\ref{fig:cavalieri}): for a Borel set $Q\subseteq\R^d$,
\[\lambda_d(Q)=\int_\R\lambda_{d-1}(Q^t)\,dt,\qquad Q^t=\{x\in\R^{d-1}:(x,t)\in Q\}.\]
For example, the cone with base a Borel set $B\subseteq\R^{d-1}$ and height $h$ has slices $(1-\frac th)B$, of measure $(1-\frac th)^{d-1}\lambda_{d-1}(B)$ (Theorem~\ref{thm:lebesgue}), so its volume is
$\lambda_{d-1}(B)\int_0^h(1-\frac th)^{d-1}dt=\frac hd\lambda_{d-1}(B)$, for any base whatsoever.

\section{Tonelli and Fubini}

\begin{theorem}[Tonelli]\label{thm:tonelli}
Let $\mu,\nu$ be $\sigma$-finite and $f:E\times F\to[0,+\infty]$ be $\mathcal A\otimes\mathcal B$-measurable. Then $x\mapsto\int_Ff(x,y)\,d\nu(y)$ and $y\mapsto\int_Ef(x,y)\,d\mu(x)$ are measurable, and
\[\int_{E\times F}f\,d(\mu\otimes\nu)=\int_E\Bigl(\int_Ff\,d\nu\Bigr)d\mu=\int_F\Bigl(\int_Ef\,d\mu\Bigr)d\nu.\]
\end{theorem}
\begin{proof}
For $f=\mathbf 1_Q$ this is Theorem~\ref{thm:product}. Both sides are additive and positively homogeneous in $f$, so it holds for non-negative simple functions. For general $f$, take simple
$0\le s_n\uparrow f$ (Proposition~\ref{prop:measurablefns}). For each $x$, $\int s_n(x,y)\,d\nu(y)\uparrow\int f(x,y)\,d\nu(y)$ by monotone convergence, so the limit is measurable in $x$; apply monotone
convergence once more in $x$, and once on $E\times F$.
\end{proof}

No integrability is needed: both sides may be $+\infty$, and then both are. In practice, Tonelli is used first, on $|f|$, to check that $f$ is integrable;
then Fubini allows the swap.

\begin{theorem}[Fubini]\label{thm:fubini}
Let $\mu,\nu$ be $\sigma$-finite and $f\in L^1(\mu\otimes\nu)$, real or complex. Then for $\mu$-almost every $x$, $f_x$ is $\nu$-integrable; the function $x\mapsto\int_Ff(x,y)\,d\nu(y)$, defined almost
everywhere, is $\mu$-integrable; the same holds with the roles exchanged; and
\[\int_{E\times F}f\,d(\mu\otimes\nu)=\int_E\Bigl(\int_Ff\,d\nu\Bigr)d\mu=\int_F\Bigl(\int_Ef\,d\mu\Bigr)d\nu.\]
\end{theorem}
\begin{proof}
By Tonelli, $\int_E(\int_F|f|\,d\nu)\,d\mu=\int|f|\,d(\mu\otimes\nu)<\infty$, so the measurable function $x\mapsto\int_F|f_x|\,d\nu$ is finite outside a $\mu$-null set $N$ (Proposition~\ref{prop:markov}). For real
$f$, write $f=f^+-f^-$; on $E\setminus N$, $\int f_x\,d\nu=\int f^+_x\,d\nu-\int f^-_x\,d\nu$, a difference of two measurable functions of $x$ that are integrable, by Tonelli applied to $f^\pm$. Integrating over
$E\setminus N$ (removing a null set changes nothing) and using Tonelli for $f^+$ and $f^-$ gives the formula. Complex $f$: apply this to the real and imaginary parts.
\end{proof}

\begin{think}
Can $\sum_m\sum_na_{mn}$ and $\sum_n\sum_ma_{mn}$ both converge, to different values? Try a matrix with $1$ on the diagonal, $-1$ just to the right of it, and $0$ elsewhere.
\end{think}

With counting measure on $\N$, Tonelli and Fubini are statements about double series, as in Tome~0, Chapter~5: a double series of non-negative terms can be summed in either order, with the
same value, finite or not; and if $\sum_{m,n}|a_{mn}|<\infty$, both iterated sums of $a_{mn}$ converge and agree. Without absolute convergence, the matrix of the box gives row sums $1-1=0$, hence
$\sum_m\sum_na_{mn}=0$, while the first column sums to $1$ and every other column to $-1+1=0$, so $\sum_n\sum_ma_{mn}=1$. Exercise~\ref{ex:fubinifail} does the same with an integral.

\begin{figure}[ht]
\centering
\fig{ch08-diagonal}
\caption{Lebesgue measure $\lambda$ on the $x$-axis, counting measure $\#$ on the $y$-axis, and the diagonal $D$ of the unit square. Every vertical section is one point, of counting measure $1$; every
horizontal section is one point, of Lebesgue measure $0$. Integrating the sections gives $1$ one way and $0$ the other.}
\label{fig:diagcount}
\end{figure}

\emph{$\sigma$-finiteness cannot be dropped.} Let $\mu=\lambda$ on $[0,1]$, let $\nu$ be counting measure on $[0,1]$, and let $D=\{(x,x)\}$ be the diagonal, a closed set, hence in $\mathcal B([0,1])\otimes\mathcal B([0,1])$ by
Proposition~\ref{prop:borelproduct}. Then $\int\nu(D_x)\,d\lambda(x)=\int1\,d\lambda=1$, but $\int\lambda(D^y)\,d\nu(y)=\sum_y0=0$ (Figure~\ref{fig:diagcount}). The proof of Lemma~\ref{lem:slicemeasurable} breaks exactly where $\nu$ was
cut into finite pieces.

\emph{Products of finitely many measures and independence.} The construction iterates: $\mu_1\otimes\cdots\otimes\mu_n$ is the unique measure on $\mathcal A_1\otimes\cdots\otimes\mathcal A_n$ with
$(\mu_1\otimes\cdots\otimes\mu_n)(A_1\times\cdots\times A_n)=\prod_i\mu_i(A_i)$, and it does not matter how the factors are grouped (both groupings agree on rectangles). In probability, random variables
$X_1,\dots,X_n$ are independent exactly when the law of $(X_1,\dots,X_n)$ is the product of their laws: both measures agree on rectangles, which is the definition of independence, and a measure on
the product is determined by its values on rectangles. Fubini then gives at once $\mathbb E[XY]=\mathbb E[X]\,\mathbb E[Y]$ for independent integrable $X$ and $Y$, since
$\int|xy|\,d(P_X\otimes P_Y)=\mathbb E|X|\,\mathbb E|Y|<\infty$, and more generally
\[\mathbb E[h(X,Y)]=\int\mathbb E[h(x,Y)]\,dP_X(x)\]
for $h\ge0$ measurable, or $h(X,Y)$ integrable: freeze $X$, average over $Y$, then average over $X$. The \emph{convolution} $\mu*\nu$ of two probabilities on $\R^d$ is the law of $X+Y$ for
independent $X\sim\mu$ and $Y\sim\nu$, that is $(\mu*\nu)(A)=\iint\mathbf 1_A(x+y)\,d\mu(x)\,d\nu(y)$; when $\mu$ and $\nu$ have densities $f$ and $g$, a translation in the inner integral and Tonelli
give the density
\[(f*g)(s)=\int_{\R^d}f(s-y)\,g(y)\,dy.\]
Indeed, $\iint\mathbf 1_A(x+y)f(x)g(y)\,dx\,dy=\int\bigl(\int\mathbf 1_A(s)f(s-y)\,ds\bigr)g(y)\,dy=\int_A\bigl(\int f(s-y)g(y)\,dy\bigr)ds$, using in the inner integral that $\lambda_d$ is translation
invariant. Chapter~9 studies convolution of functions in depth.

\section{Lebesgue measure under linear maps}

\begin{theorem}\label{thm:lineardet}
Let $T:\R^d\to\R^d$ be linear. If $T$ is invertible, then $\lambda_d(T(A))=|\det T|\,\lambda_d(A)$ for every Borel set $A$. If $T$ is not invertible, $T(\R^d)$ is a null set.
\end{theorem}
\begin{proof}
Let $T$ be invertible. Since $T$ is a homeomorphism, $T(A)=(T^{-1})^{-1}(A)$ is Borel, and $\nu_T(A)=\lambda_d(T(A))$ is a Borel measure. It is translation invariant, as $T(A+a)=T(A)+Ta$, and
$c(T)=\nu_T(\,]0,1]^d)$ is finite (the image of a bounded set is bounded) and positive (the image contains the non-empty open set $T(\,]0,1[^d)$). By the uniqueness in Theorem~\ref{thm:lebesgue},
$\nu_T=c(T)\lambda_d$. Since $\nu_{ST}(A)=\lambda_d(S(T(A)))=c(S)c(T)\lambda_d(A)$, we have $c(ST)=c(S)c(T)$; as $\det$ is multiplicative too, it suffices to prove $c(T)=|\det T|$ when $T$ is one of the
\emph{elementary} maps below, because every invertible matrix is a product of elementary ones (Gaussian elimination: each row operation is a left multiplication by an elementary matrix,
and invertible $T$ is reduced to the identity by row operations).

\emph{(i) Scaling one coordinate}, $x_i\mapsto ax_i$ with $a\ne0$: the unit cube goes to a box with one side $|a|$ and the others $1$, so $c=|a|=|\det|$.
\emph{(ii) Swapping two coordinates}: the unit cube goes to itself, $c=1=|\det|$.
\emph{(iii) A shear} $x_i\mapsto x_i+ax_j$ ($i\ne j$), other coordinates fixed, of determinant $1$. Order coordinates so that $i=1$, $j=2$. The image of $]0,1]^d$ is the set of $y$ with
$(y_2,\dots,y_d)\in\,]0,1]^{d-1}$ and $ay_2<y_1\le ay_2+1$. Each slice in the $y_1$ direction is an interval of length $1$, so by Cavalieri ($\lambda_d=\lambda\otimes\lambda_{d-1}$), $c=\int_{]0,1]^{d-1}}1\,d\lambda_{d-1}=1$.

If $T$ is not invertible, its image lies in a hyperplane $H$. Some invertible $S$ maps the coordinate hyperplane $H_0=\{x_1=0\}$ onto $H$ (send a basis of $H_0$ to a basis of $H$, and $e_1$ outside $H$), and
$\lambda_d(H_0)=0$ (it is a product $\{0\}\times\R^{d-1}$), so $\lambda_d(H)=|\det S|\cdot0=0$.
\end{proof}

So rotations and reflections preserve Lebesgue measure ($|\det|=1$), as Chapter~6 promised, and the parallelepiped spanned by $v_1,\dots,v_d$ has volume $|\det(v_1,\dots,v_d)|$: this is the
geometric meaning of the determinant: up to sign, it is the factor by which $T$ multiplies volumes.

\section{The change of variables formula}

Let $U,V\subseteq\R^d$ be open. A \emph{$C^1$ diffeomorphism} $\varphi:U\to V$ is a bijection such that $\varphi$ and $\varphi^{-1}$ are $C^1$. Its \emph{Jacobian} is $J_\varphi(x)=\det D\varphi(x)$, never zero, since
$D\varphi^{-1}(\varphi(x))\,D\varphi(x)=I$ by the chain rule. Near $x$, $\varphi$ looks like the affine map $x+h\mapsto\varphi(x)+D\varphi(x)h$, which multiplies volumes by $|J_\varphi(x)|$
(Theorem~\ref{thm:lineardet}). Summing over small cubes suggests the formula (Figure~\ref{fig:diffeo}).

\begin{figure}[ht]
\centering
\fig{ch08-diffeo}
\caption{A grid of small squares and its image under a diffeomorphism $\varphi$. The image of one small square (orange) is almost the parallelogram spanned by the columns of
$D\varphi(x)$ scaled by the side (dashed), of area $|J_\varphi(x)|$ times the area of the square.}
\label{fig:diffeo}
\end{figure}

\begin{theorem}[change of variables]\label{thm:cov}
Let $\varphi:U\to V$ be a $C^1$ diffeomorphism between open subsets of $\R^d$. For every Borel $f:V\to[0,+\infty]$,
\[\int_Vf(y)\,dy=\int_Uf(\varphi(x))\,|J_\varphi(x)|\,dx.\]
The same holds for every $f\in L^1(V)$, and then $(f\circ\varphi)\,|J_\varphi|\in L^1(U)$. In particular $\lambda_d(\varphi(A))=\int_A|J_\varphi|\,d\lambda_d$ for every Borel $A\subseteq U$.
\end{theorem}

We use the norm $\|x\|=\max_i|x_i|$ on $\R^d$, whose closed balls are cubes: the cube of centre $a$ and side $2h$ is $\{x:\|x-a\|\le h\}$, of volume $(2h)^d$. For a matrix $M$, let
$\|M\|=\max_i\sum_j|M_{ij}|$, so that $\|Mx\|\le\|M\|\,\|x\|$ and $\|I\|=1$.

\begin{lemma}\label{lem:cube}
Let $Q\subseteq U$ be a closed cube of centre $a$ and side $2h$, and $T$ an invertible matrix. Then
\[\lambda_d(\varphi(Q))\le|\det T|\Bigl(\sup_{z\in Q}\|T^{-1}D\varphi(z)\|\Bigr)^d\lambda_d(Q).\]
\end{lemma}
\begin{proof}
Let $g=T^{-1}\circ\varphi$ and $M=\sup_{z\in Q}\|Dg(z)\|$, finite since $Dg=T^{-1}D\varphi$ is continuous on the compact $Q$. For $x\in Q$, apply the mean value theorem (Theorem~\ref{thm:mvt}) to
$t\mapsto g_i(a+t(x-a))$ on $[0,1]$: for some $z$ on the segment $[a,x]\subseteq Q$,
\[|g_i(x)-g_i(a)|=\Bigl|\sum_j\partial_jg_i(z)(x_j-a_j)\Bigr|\le M\|x-a\|\le Mh.\]
So $g(Q)$ lies in the cube of centre $g(a)$ and side $2Mh$, of volume $M^d\lambda_d(Q)$, and $\lambda_d(\varphi(Q))=\lambda_d(T(g(Q)))=|\det T|\,\lambda_d(g(Q))$ by Theorem~\ref{thm:lineardet}.
\end{proof}

\begin{proof}[Proof of Theorem~\ref{thm:cov}]
\emph{Step 1: closed cubes.} Let $Q\subseteq U$ be a closed cube. We show $\lambda_d(\varphi(Q))\le\int_Q|J_\varphi|$. On the compact $Q$, $D\varphi$ is uniformly continuous and $\|D\varphi(x)^{-1}\|\le C$ for some $C$
(it is continuous). Given $\varepsilon>0$, choose $\delta$ with $\|D\varphi(z)-D\varphi(x)\|\le\varepsilon$ when $\|z-x\|\le\delta$, and cut $Q$ into $N^d$ closed subcubes $Q_k$ of centres $x_k$ and sides
at most $\delta$. For $z\in Q_k$,
\[\|D\varphi(x_k)^{-1}D\varphi(z)\|\le\|I\|+\|D\varphi(x_k)^{-1}\|\,\|D\varphi(z)-D\varphi(x_k)\|\le1+C\varepsilon.\]
Lemma~\ref{lem:cube} with $T=D\varphi(x_k)$ and subadditivity give
\[\lambda_d(\varphi(Q))\le\sum_k\lambda_d(\varphi(Q_k))\le(1+C\varepsilon)^d\sum_k|J_\varphi(x_k)|\,\lambda_d(Q_k).\]
The last sum is a Riemann sum of the continuous function $|J_\varphi|$ on $Q$, within $\omega\,\lambda_d(Q)$ of $\int_Q|J_\varphi|$, where $\omega$ is the oscillation of $|J_\varphi|$ on cubes of side $\delta$ (the faces of the
$Q_k$ are null, by Theorem~\ref{thm:lineardet}). Letting $\varepsilon\to0$ and $\delta\to0$ gives the claim.

\emph{Step 2: Borel sets.} Every open $W\subseteq U$ is a countable union of closed dyadic cubes $Q_k\subseteq W$ with disjoint interiors (take the closed dyadic cubes of side at most $1$ contained in $W$ that are maximal for inclusion among them). Their
boundaries are null, so by Step~1
\[\lambda_d(\varphi(W))\le\sum_k\lambda_d(\varphi(Q_k))\le\sum_k\int_{Q_k}|J_\varphi|=\int_W|J_\varphi|.\]
Now write $U=\bigcup_nU_n$ with $U_n$ open, bounded and $\overline{U_n}\subseteq U$ compact (for instance $U_n=\{x\in U:\|x\|<n,\ \mathrm{dist}(x,\R^d\setminus U)>\frac1n\}$), so that $|J_\varphi|$ is bounded on each
$U_n$. Let $A\subseteq U_n$ be Borel. By outer regularity (Proposition~\ref{prop:regularity}) there are open sets $W_m$ with $A\subseteq W_m\subseteq U_n$ and $\lambda_d(W_m\setminus A)\to0$; we may assume them decreasing. Then
$\lambda_d(\varphi(A))\le\lambda_d(\varphi(W_m))\le\int_{W_m}|J_\varphi|$, which tends to $\int_A|J_\varphi|$ by dominated convergence ($|J_\varphi|\mathbf 1_{U_n}$ is integrable). For a general Borel $A\subseteq U$, cut it into the
disjoint pieces $A\cap(U_n\setminus U_{n-1})$ and add up: $\lambda_d(\varphi(A))\le\int_A|J_\varphi|$.

\emph{Step 3: the inequality for functions.} For a Borel set $B\subseteq V$, Step~2 applied to $A=\varphi^{-1}(B)$ gives
$\int_V\mathbf 1_B=\lambda_d(B)\le\int_U\mathbf 1_B(\varphi(x))\,|J_\varphi(x)|\,dx$. By linearity this holds for non-negative simple functions, and by monotone convergence for every Borel $f\ge0$:
\[\int_Vf\le\int_U(f\circ\varphi)\,|J_\varphi|.\tag{$*$}\]
\emph{Step 4: equality.} Apply $(*)$ to the diffeomorphism $\varphi^{-1}:V\to U$ and the function $g=(f\circ\varphi)\,|J_\varphi|$ on $U$:
\[\int_Ug\le\int_V(g\circ\varphi^{-1})\,|J_{\varphi^{-1}}|=\int_Vf(y)\,|J_\varphi(\varphi^{-1}(y))|\,|J_{\varphi^{-1}}(y)|\,dy=\int_Vf,\]
since $D\varphi(\varphi^{-1}(y))\,D\varphi^{-1}(y)=I$. With $(*)$, this is the formula for $f\ge0$. For integrable $f$, apply it to $|f|$ (so $(f\circ\varphi)|J_\varphi|$ is integrable), then to the positive and
negative parts of the real and imaginary parts.
\end{proof}

Two remarks. The formula needs $\varphi$ to be a bijection only up to null sets: if $U\setminus U'$ and $V\setminus V'$ are null and $\varphi:U'\to V'$ is a diffeomorphism, the integrals over $U$ and $U'$, $V$ and $V'$
agree. And in dimension one, for $\varphi$ increasing, it is the substitution rule of Tome~0, now valid for all non-negative or integrable Borel functions.

\section{Polar coordinates, the \texorpdfstring{$n$}{n}-ball and Gaussian vectors}

\emph{The plane.} The map $(r,\theta)\mapsto(r\cos\theta,r\sin\theta)$ is a diffeomorphism from $]0,\infty[\times\,]0,2\pi[$ onto the plane minus the half-line $[0,\infty[\times\{0\}$, a null set, with Jacobian
$\cos\theta\cdot r\cos\theta-(-r\sin\theta)\sin\theta=r$. Hence $\int_{\R^2}f=\int_0^{2\pi}\int_0^\infty f(r\cos\theta,r\sin\theta)\,r\,dr\,d\theta$ for $f\ge0$. With $f(x,y)=e^{-x^2-y^2}$ and Tonelli,
\[\Bigl(\int_\R e^{-x^2}dx\Bigr)^2=\int_0^{2\pi}\!\!\int_0^\infty e^{-r^2}r\,dr\,d\theta=2\pi\cdot\frac12=\pi,\]
so $\int_\R e^{-x^2}dx=\sqrt\pi$: Poisson's computation, used throughout the earlier tomes, is now a theorem.

In higher dimension one could write spherical coordinates explicitly (Exercise~\ref{ex:spherical3} does it in $\R^3$), but a cleaner route separates the radius from the direction. Let
$S=S^{d-1}=\{x\in\R^d:|x|=1\}$ with the Euclidean norm, and $\Phi(x)=(|x|,\frac x{|x|})$, a homeomorphism from $\R^d\setminus\{0\}$ onto $]0,\infty[\times S$, with inverse $(r,\omega)\mapsto r\omega$.

\begin{theorem}[polar coordinates]\label{thm:polar}
There is a unique Borel measure $\sigma$ on $S^{d-1}$ such that for every Borel $f:\R^d\to[0,+\infty]$,
\[\int_{\R^d}f(x)\,dx=\int_0^\infty\!\!\int_{S^{d-1}}f(r\omega)\,d\sigma(\omega)\,r^{d-1}\,dr.\]
It is given by $\sigma(B)=d\cdot\lambda_d(\{r\omega:0<r\le1,\ \omega\in B\})$, and it is invariant under rotations.
\end{theorem}
\begin{proof}
Define $\sigma$ by the formula; it is a finite Borel measure on $S$ since $B\mapsto\{r\omega:0<r\le1,\omega\in B\}=\Phi^{-1}(]0,1]\times B)$ preserves disjoint unions. Let $\rho$ be the measure $r^{d-1}dr$ on $]0,\infty[$, and
compare the image $\Phi_*\lambda_d$ (that is, $A\mapsto\lambda_d(\Phi^{-1}(A))$) with $\rho\otimes\sigma$. On a rectangle $]a,b]\times B$, the set $\Phi^{-1}(]a,b]\times B)$ is the difference of the dilations by $b$ and by $a$ of
$\Phi^{-1}(]0,1]\times B)$, so by the scaling of Theorem~\ref{thm:lebesgue}
\[\lambda_d\bigl(\Phi^{-1}(]a,b]\times B)\bigr)=(b^d-a^d)\frac{\sigma(B)}d=\rho(]a,b])\,\sigma(B).\]
These rectangles form a $\pi$-system that generates the Borel sets of $]0,\infty[\times S$ (as in Proposition~\ref{prop:borelproduct}: every open set is a countable union of products of basic open sets, both
factors being separable metric spaces), and both measures are finite on $]\frac1n,n]\times S$. By Theorem~\ref{thm:uniq6}, $\Phi_*\lambda_d=\rho\otimes\sigma$ on $]0,\infty[\times S$. Since $\{0\}$ is null,
$\int_{\R^d}f\,d\lambda_d=\int(f\circ\Phi^{-1})\,d(\rho\otimes\sigma)$, and Tonelli gives the formula. Uniqueness: taking $f=\mathbf 1_{\Phi^{-1}(]0,1]\times B)}$ in the formula forces $\sigma(B)=d\cdot\lambda_d(\ldots)$. A rotation $R$
preserves $\lambda_d$ (Theorem~\ref{thm:lineardet}) and maps $\Phi^{-1}(]0,1]\times B)$ to $\Phi^{-1}(]0,1]\times RB)$, so $\sigma(RB)=\sigma(B)$.
\end{proof}

For a radial function,
\[\int_{\R^d}g(|x|)\,dx=\sigma(S^{d-1})\int_0^\infty g(r)\,r^{d-1}\,dr.\]
Applying this to $e^{-|x|^2}$, whose integral is $(\sqrt\pi)^d$ by Tonelli, and substituting $t=r^2$ in the Gamma
integral $\Gamma(s)=\int_0^\infty t^{s-1}e^{-t}\,dt$:
\[\pi^{d/2}=\sigma(S^{d-1})\int_0^\infty e^{-r^2}r^{d-1}\,dr=\sigma(S^{d-1})\,\frac{\Gamma(\frac d2)}2.\]

\begin{theorem}[the $n$-ball]\label{thm:ball}
The unit sphere and the unit ball of $\R^d$ have measures
\[\sigma(S^{d-1})=\frac{2\pi^{d/2}}{\Gamma(\frac d2)},\qquad\lambda_d(B^d)=\frac{\sigma(S^{d-1})}d=\frac{\pi^{d/2}}{\Gamma(\frac d2+1)}.\]
\end{theorem}
\begin{proof}
The first formula was just computed. For the second, $\lambda_d(B^d)=\sigma(S^{d-1})\int_0^1r^{d-1}\,dr$, and $\frac d2\Gamma(\frac d2)=\Gamma(\frac d2+1)$.
\end{proof}

This gives $2$, $\pi$, $\frac43\pi$, $\frac{\pi^2}2$ in dimensions $1$ to $4$, the formula of Tome~4, Chapter~5, now with every step proved. Since $\sigma(S^1)=2\pi$ is the length of the circle and
$\sigma(S^2)=4\pi$ the area of the sphere, $\sigma$ is the surface measure.

\begin{think}
What is the volume of the unit ball in dimension $100$? Compare it with the volume of the cube $[-1,1]^{100}$ that contains it. Where in the cube is its volume?
\end{think}

$\lambda_{100}(B^{100})=\frac{\pi^{50}}{50!}\approx2.4\cdot10^{-40}$, while the cube has volume $2^{100}\approx1.3\cdot10^{30}$: the ball fills a fraction $2\cdot10^{-70}$ of it. Almost all of the cube's volume is in its
corners, at distance about $\sqrt{100/3}\approx5.8$ from the centre (the average of $x_1^2+\cdots+x_{100}^2$ is $\frac{100}3$). And within the ball, almost all the volume is near the surface: the ball of
radius $0.9$ has a fraction $0.9^{100}\approx3\cdot10^{-5}$ of the volume. High-dimensional geometry defies the intuition of dimension three; this is the \emph{curse of dimensionality}
that Monte Carlo methods escape (notebook of Chapter~7) and that rejection sampling falls into (this chapter's notebook).

\emph{Gaussian vectors.} A standard Gaussian vector $Z=(Z_1,\dots,Z_d)$ has independent $\mathcal N(0,1)$ coordinates; its law is the product of the laws, with density
$\prod_i\frac{e^{-z_i^2/2}}{\sqrt{2\pi}}=(2\pi)^{-d/2}e^{-|z|^2/2}$, a function of $|z|$ only. A \emph{Gaussian vector} is $X=AZ+m$ with $A$ a $d\times d$ matrix and $m\in\R^d$; its covariance matrix is
$\Sigma=AA^{\mathsf T}$.

\begin{theorem}[Gaussian densities]\label{thm:gaussdensity}
If $\Sigma=AA^{\mathsf T}$ is invertible, $X=AZ+m$ has the density
\[f_X(x)=\frac1{(2\pi)^{d/2}\sqrt{\det\Sigma}}\exp\Bigl(-\frac12(x-m)^{\mathsf T}\Sigma^{-1}(x-m)\Bigr).\]
In particular the law of $X$ depends only on $m$ and $\Sigma$, and if $R$ is orthogonal, $RZ$ is again standard Gaussian.
\end{theorem}
\begin{proof}
$\det\Sigma=(\det A)^2\ne0$, so $A$ is invertible and $\varphi(z)=Az+m$ is a diffeomorphism of $\R^d$ with $|J_\varphi|=|\det A|=\sqrt{\det\Sigma}$. For a Borel set $B$, Theorem~\ref{thm:cov} applied to $\varphi^{-1}$ gives
\[\mathbb P(X\in B)=\int\mathbf 1_B(Az+m)\,\frac{e^{-|z|^2/2}}{(2\pi)^{d/2}}\,dz=\int_B\frac{e^{-|A^{-1}(x-m)|^2/2}}{(2\pi)^{d/2}|\det A|}\,dx,\]
and $|A^{-1}(x-m)|^2=(x-m)^{\mathsf T}(A^{-1})^{\mathsf T}A^{-1}(x-m)=(x-m)^{\mathsf T}\Sigma^{-1}(x-m)$. For orthogonal $R$, $RR^{\mathsf T}=I$.
\end{proof}

This is the density Tome~3, Chapter~6, used for Gaussian vectors, now proved. The rotation invariance of $Z$ is the reason behind the uniform law on the sphere (Exercise~\ref{ex:sphere}) and behind
the Box--Muller method for simulating Gaussians (notebook).

\begin{beauty}
One formula, $\int_Vf=\int_U(f\circ\varphi)\,|J_\varphi|$, contains the substitution rule of Tome~0, polar coordinates, the volume of balls in every dimension, the Gaussian integral, the density of every
Gaussian vector and the geometric meaning of the determinant. Behind it are two simple facts: Lebesgue measure is the only translation-invariant way to measure volume, and a smooth map is, on a
small enough scale, linear.
\end{beauty}

\section{Exercises}

\begin{exercise}[$\star$]\label{ex:fubinifail}
Let $f(x,y)=\frac{x^2-y^2}{(x^2+y^2)^2}$ on $]0,1[^2$. Show that $\int_0^1(\int_0^1f\,dy)\,dx=\frac\pi4$ and $\int_0^1(\int_0^1f\,dx)\,dy=-\frac\pi4$. Which hypothesis of Fubini's theorem fails? Check it with polar
coordinates.
\end{exercise}

\begin{exercise}[$\star\star$]\label{ex:basel}
\emph{Ritelli's proof of $\zeta(2)=\frac{\pi^2}6$ (Tome~4, Chapter~3).} Compute $I=\int_0^\infty\!\int_0^\infty\frac{dx\,dy}{(1+y)(1+x^2y)}$ in both orders. Deduce $\int_0^1\frac{-\ln x}{1-x^2}\,dx=\frac{\pi^2}8$, then
$\sum_{k\ge0}\frac1{(2k+1)^2}=\frac{\pi^2}8$ and $\zeta(2)=\frac{\pi^2}6$. Justify every exchange of limits.
\end{exercise}

\begin{exercise}[$\star\star$]\label{ex:betagamma}
For $a,b>0$ let $B(a,b)=\int_0^1v^{a-1}(1-v)^{b-1}\,dv$. Show that $\Gamma(a)\Gamma(b)=B(a,b)\,\Gamma(a+b)$, writing $\Gamma(a)\Gamma(b)$ as an integral over $]0,\infty[^2$ and changing variables by
$(s,t)\mapsto(u,v)=(s+t,\frac s{s+t})$.
\end{exercise}

\begin{exercise}[$\star$]\label{ex:gammasum}
The Gamma law of parameter $a>0$ has density $\frac{t^{a-1}e^{-t}}{\Gamma(a)}$ on $]0,\infty[$. Show that the sum of independent Gamma variables of parameters $a$ and $b$ is Gamma of parameter $a+b$. Deduce
the law of a sum of $n$ independent exponential variables of parameter $1$.
\end{exercise}

\begin{exercise}[$\star\star$]\label{ex:sphere}
\emph{The uniform law on the sphere.} Let $Z$ be a standard Gaussian vector of $\R^d$. Show that $|Z|$ and $\frac Z{|Z|}$ are independent, and that $\frac Z{|Z|}$ has law $\frac\sigma{\sigma(S^{d-1})}$, the normalised
surface measure. If moreover $U$ is uniform on $]0,1[$ and independent of $Z$, show that $U^{1/d}\frac Z{|Z|}$ is uniform in the unit ball.
\end{exercise}

\begin{exercise}[$\star$]\label{ex:spherical3}
Show that $(r,\theta,\phi)\mapsto(r\sin\theta\cos\phi,r\sin\theta\sin\phi,r\cos\theta)$ is a diffeomorphism from $]0,\infty[\times\,]0,\pi[\times\,]0,2\pi[$ onto $\R^3$ minus a null set, with Jacobian
$r^2\sin\theta$. Recover the volume of the ball and the area of the sphere, and compute $\int_{\R^3}\frac{e^{-|x|}}{|x|}\,dx$.
\end{exercise}

\begin{exercise}[$\star\star$]\label{ex:shell}
\emph{Concentration of volume.} (a) Show that the fraction of the volume of $B^d$ at distance less than $\varepsilon$ from the sphere tends to $1$ as $d\to\infty$, for every fixed $\varepsilon>0$. (b) Show that
the fraction of the volume of $B^d$ in the slab $|x_1|\le\frac c{\sqrt d}$ tends to $\int_{-c}^ce^{-t^2/2}\frac{dt}{\sqrt{2\pi}}$. So almost all the volume of a high-dimensional ball is near its surface, and near
every equator.
\end{exercise}

\begin{exercise}[$\star\star\star$]\label{ex:simplex}
\emph{Dirichlet's integral.} Let $\Delta_n=\{x\in\,]0,\infty[^n:x_1+\cdots+x_n<1\}$ and $a_1,\dots,a_{n+1}>0$. Show that
\[\int_{\Delta_n}x_1^{a_1-1}\cdots x_n^{a_n-1}(1-x_1-\cdots-x_n)^{a_{n+1}-1}\,dx=\frac{\Gamma(a_1)\cdots\Gamma(a_{n+1})}{\Gamma(a_1+\cdots+a_{n+1})},\]
by computing $\int_{]0,\infty[^{n+1}}\prod_it_i^{a_i-1}e^{-t_i}\,dt$ in the variables $s=t_1+\cdots+t_{n+1}$, $x_i=\frac{t_i}s$ ($i\le n$). Deduce $\lambda_n(\Delta_n)=\frac1{n!}$.
\end{exercise}

\section{Hints}

\begin{hint}{ex:fubinifail}
$\partial_y\frac y{x^2+y^2}=\frac{x^2-y^2}{(x^2+y^2)^2}$, and $f(y,x)=-f(x,y)$. In polar coordinates $|f|=\frac{|\cos2\theta|}{r^2}$.
\end{hint}
\begin{hint}{ex:basel}
In $x$ first, substitute $x=u/\sqrt y$. In $y$ first, use partial fractions in $y$. For $\int_1^\infty$, substitute $x=\frac1u$. Expand $\frac1{1-x^2}$ as a geometric series and integrate term by term,
using $\int_0^1x^{2k}(-\ln x)\,dx=\frac1{(2k+1)^2}$. Separate odd and even $n$ in $\zeta(2)$.
\end{hint}
\begin{hint}{ex:betagamma}
The inverse map is $(u,v)\mapsto(uv,u(1-v))$, from $]0,\infty[\times\,]0,1[$ onto $]0,\infty[^2$, with Jacobian of absolute value $u$.
\end{hint}
\begin{hint}{ex:gammasum}
Convolve the densities; in $\int_0^sx^{a-1}(s-x)^{b-1}dx$ substitute $x=sv$. Use Exercise~\ref{ex:betagamma}.
\end{hint}
\begin{hint}{ex:sphere}
For bounded Borel $g,h\ge0$, compute $\mathbb E[g(|Z|)h(\frac Z{|Z|})]$ with Theorem~\ref{thm:polar}. For the ball, compute $\mathbb P(|X|\le r,\ \frac X{|X|}\in B)$ and compare with
$\lambda_d(\{s\omega:s\le r,\omega\in B\})/\lambda_d(B^d)$.
\end{hint}
\begin{hint}{ex:spherical3}
Expand the $3\times3$ determinant along the last row. The map fails to be onto only on a half-plane, which is null. For the integral, Theorem~\ref{thm:polar} is quicker.
\end{hint}
\begin{hint}{ex:shell}
(a) The ball of radius $1-\varepsilon$ is a dilate. (b) By Cavalieri, the slice of $B^d$ at height $x_1=t$ is a ball of dimension $d-1$ and radius $\sqrt{1-t^2}$. Substitute $t=\frac s{\sqrt d}$
and use dominated convergence with $(1-\frac{s^2}d)^{(d-1)/2}\le e^{-s^2(d-1)/(2d)}$; the ratio of ball volumes is computed from Theorem~\ref{thm:ball}.
\end{hint}
\begin{hint}{ex:simplex}
The map $t\mapsto(s,x)$ is a diffeomorphism from $]0,\infty[^{n+1}$ onto $]0,\infty[\times\Delta_n$, with inverse $t_i=sx_i$ ($i\le n$), $t_{n+1}=s(1-\sum x_i)$. Add the first $n$ rows of the Jacobian matrix of
the inverse to the last one: its determinant is $s^n$.
\end{hint}

\section{Solutions}

\begin{solution}{ex:fubinifail}
\idea The function is a derivative in $y$, so the inner integrals are explicit; antisymmetry then flips the sign of the other iterated integral. Two different answers can only mean that $f$ is
not integrable, and the singularity at the origin is visible in polar coordinates.
\writeup For fixed $x\in\,]0,1[$, $y\mapsto\frac y{x^2+y^2}$ has derivative $\frac{(x^2+y^2)-2y^2}{(x^2+y^2)^2}=f(x,y)$, so $\int_0^1f(x,y)\,dy=\frac1{1+x^2}$ and $\int_0^1\frac{dx}{1+x^2}=\arctan1=\frac\pi4$. Since
$f(y,x)=-f(x,y)$, the other iterated integral is $-\frac\pi4$. By Fubini, $f$ is not integrable on $]0,1[^2$. Indeed, $]0,1[^2$ contains the quarter disc $\{r<1,\ 0<\theta<\frac\pi2\}$, on which
$|f|=\frac{r^2|\cos2\theta|}{r^4}$, and
\[\int_0^{\pi/2}\!\!\int_0^1\frac{|\cos2\theta|}{r^2}\,r\,dr\,d\theta=\int_0^{\pi/2}|\cos2\theta|\,d\theta\cdot\int_0^1\frac{dr}r=\infty.\]
The positive part (near the $x$-axis) and the negative part (near the $y$-axis) both have infinite integral, and each order of integration subtracts infinities in a different way.
\end{solution}

\begin{solution}{ex:basel}
\idea The integrand is positive, so Tonelli allows both orders without any check: the same number comes out as $\frac{\pi^2}2$ one way and as a series the other way. Every exchange of
sum and integral below is of non-negative terms.
\writeup The integrand is continuous and positive on $]0,\infty[^2$, so by Tonelli both iterated integrals equal $I$, finite or not. \emph{$x$ first.} For $y>0$, substituting $x=u/\sqrt y$,
$\int_0^\infty\frac{dx}{1+x^2y}=\frac1{\sqrt y}\cdot\frac\pi2$, and then with $y=w^2$,
\[I=\frac\pi2\int_0^\infty\frac{dy}{\sqrt y(1+y)}=\frac\pi2\int_0^\infty\frac{2\,dw}{1+w^2}=\frac{\pi^2}2.\]
\emph{$y$ first.} For $x>0$, $x\ne1$, $\frac1{(1+y)(1+x^2y)}=\frac1{1-x^2}\bigl(\frac1{1+y}-\frac{x^2}{1+x^2y}\bigr)$, so
\[\int_0^\infty\frac{dy}{(1+y)(1+x^2y)}=\frac1{1-x^2}\Bigl[\ln\frac{1+y}{1+x^2y}\Bigr]_0^\infty=\frac{-2\ln x}{1-x^2}.\]
Hence $\frac{\pi^2}2=\int_0^\infty\frac{-2\ln x}{1-x^2}\,dx$ (the point $x=1$ is null). On $]1,\infty[$ substitute $x=\frac1u$, $dx=\frac{du}{u^2}$: the integrand becomes $\frac{2\ln u}{1-u^{-2}}\cdot\frac1{u^2}=\frac{-2\ln u}{1-u^2}$, so both
halves are equal and $\int_0^1\frac{-\ln x}{1-x^2}\,dx=\frac{\pi^2}8$. On $]0,1[$, $\frac{-\ln x}{1-x^2}=\sum_{k\ge0}x^{2k}(-\ln x)$, a series of non-negative functions, which may be integrated term by term (monotone
convergence). Integrating by parts, $\int_0^1x^{2k}(-\ln x)\,dx=\bigl[\frac{x^{2k+1}}{2k+1}(-\ln x)\bigr]_0^1+\int_0^1\frac{x^{2k}}{2k+1}\,dx=\frac1{(2k+1)^2}$. So $\sum_k\frac1{(2k+1)^2}=\frac{\pi^2}8$. Finally, splitting
$\zeta(2)$ into odd and even terms, $\zeta(2)=\frac{\pi^2}8+\frac14\zeta(2)$, so $\zeta(2)=\frac43\cdot\frac{\pi^2}8=\frac{\pi^2}6$.
\end{solution}

\begin{solution}{ex:betagamma}
\idea In the double integral for $\Gamma(a)\Gamma(b)$, the exponential $e^{-(s+t)}$ only sees $s+t$. Taking $u=s+t$ as a variable separates the integral into a Gamma integral in $u$ and a Beta integral
in the proportion $v=\frac s{s+t}$.
\writeup By Tonelli,
\[\Gamma(a)\Gamma(b)=\int_{]0,\infty[^2}s^{a-1}t^{b-1}e^{-s-t}\,ds\,dt.\]
The map $\psi(u,v)=(uv,u(1-v))$ is a $C^1$ bijection from $]0,\infty[\times\,]0,1[$ onto $]0,\infty[^2$, with $C^1$ inverse
$(s,t)\mapsto(s+t,\frac s{s+t})$, and
\[J_\psi=\det\begin{pmatrix}v&u\\1-v&-u\end{pmatrix}=-uv-u(1-v)=-u.\]
By Theorem~\ref{thm:cov} and Tonelli,
\[\begin{aligned}\Gamma(a)\Gamma(b)&=\int_0^\infty\!\!\int_0^1(uv)^{a-1}(u(1-v))^{b-1}e^{-u}\,u\,dv\,du\\&=\int_0^\infty u^{a+b-1}e^{-u}\,du\eqbr\cdot\int_0^1v^{a-1}(1-v)^{b-1}\,dv\\&=\Gamma(a+b)B(a,b).\end{aligned}\]
\end{solution}

\begin{solution}{ex:gammasum}
\idea The density of a sum of independent variables is the convolution of their densities; for Gamma densities the convolution integral is a Beta integral, and Exercise~\ref{ex:betagamma}
supplies its value.
\writeup Let $f_a(t)=\frac{t^{a-1}e^{-t}}{\Gamma(a)}\mathbf 1_{t>0}$. The sum has density $(f_a*f_b)(s)=\int f_a(x)f_b(s-x)\,dx$, which is $0$ for $s\le0$ and, for $s>0$,
\[\frac{e^{-s}}{\Gamma(a)\Gamma(b)}\int_0^sx^{a-1}(s-x)^{b-1}\,dx=\frac{e^{-s}s^{a+b-1}}{\Gamma(a)\Gamma(b)}\int_0^1v^{a-1}(1-v)^{b-1}\,dv\]
(substituting $x=sv$), which is $f_{a+b}(s)$ by Exercise~\ref{ex:betagamma}. An exponential variable of parameter $1$ is Gamma of parameter $1$, so by induction a sum of $n$ independent ones has
density $\frac{s^{n-1}e^{-s}}{(n-1)!}$ on $]0,\infty[$: the waiting time for the $n$-th event of a Poisson process (Tome~3).
\end{solution}

\begin{solution}{ex:sphere}
\idea The Gaussian density depends only on $|z|$, so in polar coordinates it is a function of $r$ times the constant $1$ in $\omega$: the radius and the direction factorise, which is independence,
and the direction is distributed like $\sigma$ itself. For the ball, the radius must be redistributed to have density proportional to $r^{d-1}$.
\writeup Let $g,h\ge0$ be bounded Borel functions on $]0,\infty[$ and on $S=S^{d-1}$. By Theorem~\ref{thm:polar},
\[\begin{aligned}\mathbb E\Bigl[g(|Z|)h\Bigl(\frac Z{|Z|}\Bigr)\Bigr]&=\int_0^\infty\!\!\int_Sg(r)h(\omega)\frac{e^{-r^2/2}}{(2\pi)^{d/2}}\,d\sigma(\omega)\,r^{d-1}dr\\&=c\int_0^\infty g(r)e^{-r^2/2}r^{d-1}\,dr\eqbr\cdot\int_Sh\,\frac{d\sigma}{\sigma(S)}\end{aligned}\]
with $c=\frac{\sigma(S)}{(2\pi)^{d/2}}$. With $h=1$, the first factor is $\mathbb E[g(|Z|)]$; with $g=1$, the left side is $\mathbb E[h(Z/|Z|)]$ and the first factor is $1$. So
$\mathbb E[g(|Z|)h(Z/|Z|)]=\mathbb E[g(|Z|)]\,\mathbb E[h(Z/|Z|)]$ (independence, taking indicators for $g,h$), and $Z/|Z|$ has law $\sigma/\sigma(S)$. Now let $X=U^{1/d}\omega$ with $\omega=Z/|Z|$ independent of
$U$. For $0<r\le1$ and Borel $B\subseteq S$,
\[\mathbb P\Bigl(|X|\le r,\ \frac X{|X|}\in B\Bigr)=\mathbb P(U\le r^d)\,\frac{\sigma(B)}{\sigma(S)}=r^d\frac{\sigma(B)}{\sigma(S)},\]
while for the uniform law on $B^d$, by Theorem~\ref{thm:polar}, the same probability is $\frac{\sigma(B)\,r^d/d}{\sigma(S)/d}$. The sets $\{s\omega:0<s\le r,\ \omega\in B\}$ correspond under $\Phi$ to
rectangles, which form a $\pi$-system generating the Borel sets of $B^d\setminus\{0\}$; the two laws agree on it, hence everywhere.
\end{solution}

\begin{solution}{ex:spherical3}
\idea A direct computation of a $3\times3$ determinant; then the radial formula of Theorem~\ref{thm:polar} with $\sigma(S^2)=4\pi$ does the integral in one line.
\writeup The map is $C^1$. Its inverse sends $x$ to $r=|x|$, $\theta=\arccos\frac{x_3}{|x|}$ and $\phi$, the angle of $(x_1,x_2)$ in $]0,2\pi[$; it is $C^1$ away from the half-plane $\{x_2=0,\ x_1\ge0\}$, which is a null
set and is exactly what the image misses. The Jacobian matrix has rows
\[\begin{gathered}(\sin\theta\cos\phi,\ r\cos\theta\cos\phi,\ -r\sin\theta\sin\phi),\\(\sin\theta\sin\phi,\ r\cos\theta\sin\phi,\ r\sin\theta\cos\phi),\\(\cos\theta,\ -r\sin\theta,\ 0).\end{gathered}\] Expanding along the last row,
\[\begin{aligned}J&=\cos\theta\bigl(r^2\cos\theta\sin\theta\cos^2\phi+r^2\cos\theta\sin\theta\sin^2\phi\bigr)\eqbr\quad+r\sin\theta\bigl(r\sin^2\theta\cos^2\phi+r\sin^2\theta\sin^2\phi\bigr)\\&=r^2\sin\theta(\cos^2\theta+\sin^2\theta)=r^2\sin\theta>0.\end{aligned}\]
So $\lambda_3(B^3)=\int_0^1r^2dr\int_0^\pi\sin\theta\,d\theta\int_0^{2\pi}d\phi=\frac13\cdot2\cdot2\pi=\frac43\pi$, and the sphere has area $\int_0^\pi\sin\theta\,d\theta\int_0^{2\pi}d\phi=4\pi=\sigma(S^2)$. Finally, by
Theorem~\ref{thm:polar}, $\int_{\R^3}\frac{e^{-|x|}}{|x|}\,dx=4\pi\int_0^\infty\frac{e^{-r}}rr^2\,dr=4\pi\,\Gamma(2)=4\pi$.
\end{solution}

\begin{solution}{ex:shell}
\idea Volume scales like $r^d$, so in high dimension a slightly smaller ball is a negligible part of the ball. For the slab, Cavalieri reduces everything to a one-dimensional integral of
$(1-t^2)^{(d-1)/2}$, which on the scale $t=s/\sqrt d$ looks like $e^{-s^2/2}$.
\writeup (a) The ball of radius $1-\varepsilon$ has volume $(1-\varepsilon)^d\lambda_d(B^d)$, so the shell has fraction $1-(1-\varepsilon)^d\to1$. (b) For $|t|<1$ the slice $\{x'\in\R^{d-1}:(t,x')\in B^d\}$ is the ball of radius
$\sqrt{1-t^2}$, so by Cavalieri the fraction is
\[p_d=\frac{\omega_{d-1}}{\omega_d}\int_{-c/\sqrt d}^{c/\sqrt d}(1-t^2)^{\frac{d-1}2}dt=\frac{\omega_{d-1}}{\sqrt d\,\omega_d}\int_{-c}^c\Bigl(1-\frac{s^2}d\Bigr)^{\frac{d-1}2}ds,\]
with $\omega_d=\lambda_d(B^d)$ and $d>c^2$. The integrand tends to $e^{-s^2/2}$ and is bounded by $1$, so the integral tends to $\int_{-c}^ce^{-s^2/2}ds$ (dominated convergence). The same computation with $c$
replaced by $\sqrt d$ (the whole ball, fraction $1$) gives $1=\frac{\omega_{d-1}}{\sqrt d\,\omega_d}\int_{-\sqrt d}^{\sqrt d}(1-\frac{s^2}d)^{\frac{d-1}2}ds$, and that integral tends to $\int_\R e^{-s^2/2}ds=\sqrt{2\pi}$ by dominated
convergence with the bound $e^{-s^2(d-1)/(2d)}\le e^{-s^2/4}$ (for $d\ge2$, using $1-u\le e^{-u}$). Hence $\frac{\omega_{d-1}}{\sqrt d\,\omega_d}\to\frac1{\sqrt{2\pi}}$ and $p_d\to\int_{-c}^ce^{-s^2/2}\frac{ds}{\sqrt{2\pi}}$.
Any hyperplane through the centre can play the role of $x_1=0$, by rotation invariance: most of the volume lies within $\frac3{\sqrt d}$ of \emph{every} equator.
\end{solution}

\begin{solution}{ex:simplex}
\idea As in Exercise~\ref{ex:betagamma}: the exponential only depends on the total $s$, and the proportions $x_i=t_i/s$ live in the simplex. The change of variables splits the integral into a Gamma
integral in $s$ and Dirichlet's integral over the simplex.
\writeup Let $\psi(s,x)=(sx_1,\dots,sx_n,s(1-x_1-\cdots-x_n))$ on $]0,\infty[\times\Delta_n$. It is a $C^1$ bijection onto $]0,\infty[^{n+1}$, with $C^1$ inverse $t\mapsto(\sum t_i,\frac{t_1}{\sum t_i},\dots,\frac{t_n}{\sum t_i})$. Order
the variables as $(x_1,\dots,x_n,s)$. The Jacobian matrix has, for $i\le n$, row $i$ equal to $s$ in column $i$ and $x_i$ in the last column; the last row is $(-s,\dots,-s,1-\sum x_i)$. Adding the first $n$
rows to the last one turns it into $(0,\dots,0,1)$, so the matrix becomes triangular and $|J_\psi|=s^n$. Let $A=a_1+\cdots+a_{n+1}$ and $p(x)=\prod_{i\le n}x_i^{a_i-1}\cdot(1-\sum_ix_i)^{a_{n+1}-1}$, so that $\prod_it_i^{a_i-1}=s^{A-n-1}p(x)$ when $t=\psi(s,x)$. By Theorem~\ref{thm:cov} and Tonelli,
\[\begin{aligned}\prod_i\Gamma(a_i)&=\int_0^\infty\!\!\int_{\Delta_n}s^{A-n-1}p(x)\,e^{-s}\,s^n\,dx\,ds\\&=\Gamma(A)\int_{\Delta_n}p(x)\,dx.\end{aligned}\]
With all $a_i=1$, $\lambda_n(\Delta_n)=\frac{\Gamma(1)^{n+1}}{\Gamma(n+1)}=\frac1{n!}$. In probability, the proportions $(\frac{t_1}s,\dots,\frac{t_{n+1}}s)$ of independent Gamma variables follow the \emph{Dirichlet law}, the
basic prior on probability vectors in Bayesian statistics.
\end{solution}

\section*{Sources and further reading}

\begin{itemize}
\item G.~Fubini, ``Sugli integrali multipli'', \emph{Rend. Accad. Lincei} 16 (1907), 608--614; L.~Tonelli, ``Sull'integrazione per parti'', \emph{Rend. Accad. Lincei} 18 (1909), 246--253.
\item G.~B.~Folland, \emph{Real Analysis}, 2nd ed., Wiley, 1999, Sections~2.5--2.7. Product measures, linear maps, the change of variables by cubes as in Lemma~\ref{lem:cube}, and polar coordinates as in
Theorem~\ref{thm:polar}.
\item W.~Rudin, \emph{Real and Complex Analysis}, 3rd ed., McGraw--Hill, 1987, Chapters~7--8. A different proof of the change of variables, through differentiation of measures.
\item D.~Ritelli, ``Another proof of $\zeta(2)=\frac{\pi^2}6$ using double integrals'', \emph{Amer. Math. Monthly} 120 (2013), 642--645.
\item K.~Ball, ``An elementary introduction to modern convex geometry'', in \emph{Flavors of Geometry}, MSRI Publications 31, Cambridge, 1997. The geometry of high-dimensional balls and cubes,
and concentration of volume.
\item G.~E.~P.~Box and M.~E.~Muller, ``A note on the generation of random normal deviates'', \emph{Ann. Math. Statist.} 29 (1958), 610--611.
\end{itemize}
